% Generated draft: requires mathematical and source review.
\documentclass[11pt]{article}
\usepackage[margin=25mm]{geometry}
\usepackage{fontspec}
\setmainfont{Latin Modern Roman}
\usepackage{amsmath,amssymb,mathtools,unicode-math}
\setmathfont{Latin Modern Math}
\usepackage{enumitem,needspace,xurl,hyperref}
\setlength{\parindent}{0pt}
\setlength{\parskip}{5pt}
\setlength{\emergencystretch}{3em}
\newcommand{\N}{\mathbb{N}}
\newcommand{\Z}{\mathbb{Z}}
\newcommand{\Q}{\mathbb{Q}}
\newcommand{\R}{\mathbb{R}}
\newcommand{\C}{\mathbb{C}}
\newcommand{\F}{\mathbb{F}}
\newcommand{\A}{\mathbb{A}}
\newcommand{\PP}{\mathbb{P}}
\newcommand{\E}{\mathbb{E}}
\newcommand{\eps}{\varepsilon}
\newcommand{\dd}{\,\mathrm{d}}
\begin{document}
\section*{UnsolvedMath ID: 800906}
\textbf{Source problem code:} catalogue entry
\subsection{Definitions and mathematical statement}
For an explicit local interpretation, let \(n\geq 1\), let \(\Omega\subset\C^n\) be open, and let \(\varphi\colon\Omega\to[-\infty,\infty)\) be plurisubharmonic. Here plurisubharmonic means upper semicontinuous, with restriction to each connected component of the intersection with a complex affine line either subharmonic or identically \(-\infty\); the usual convention excludes functions identically \(-\infty\) on a connected component of \(\Omega\).

Using the customary second-order interpretation, a \emph{differential test from above} for \(\varphi\) at \(a\in\Omega\) is a real-valued function \(q\in C^2(U,\R)\), where \(U\subset\Omega\) is an open neighborhood of \(a\), such that
\[
q(a)=\varphi(a),\qquad \varphi(z)\leq q(z)\quad(z\in U).
\]
Define the exceptional set
\[
E^+(\varphi)=\{a\in\Omega:\text{there is no differential test from above for }\varphi\text{ at }a\}.
\]
\textbf{Question 37 (Guedj--Zeriahi).} What can be said about \(E^+(\varphi)\) for a plurisubharmonic function \(\varphi\)? Equivalently, determine what restrictions plurisubharmonicity imposes on the set of points admitting no upper differential test. The supplied question does not specify a particular type of characterization or a proposed conclusion.

\textbf{Unresolved formulation qualifications.} The source record does not specify the ambient domain, whether \(\varphi\) must be finite-valued or continuous, or the differentiability class required of a test. The open-domain formulation and \(C^2\) tests above make a standard local interpretation explicit, but require confirmation against the intended source conventions. Only interior points are covered by this formulation; no boundary condition is supplied.

The supplied background gives \(\varphi(z)=\lVert z\rVert\) on \(\C^n\) as an example having no upper differential test at the origin. It also describes this function as strictly plurisubharmonic in the viscosity sense but not in the classical sense. These descriptions are retained as source qualifications: the definitions of the two strictness notions are not supplied, so their intended comparison remains unresolved. The background additionally refers to an earlier example having no upper differential tests along \((w=0)\); that example is absent from the supplied record.
\subsection{Short English statement}
Question 37, attributed to Guedj--Zeriahi, asks what can be said about the points where a plurisubharmonic function cannot be touched from above by a smooth test function. The draft uses a local \(C^2\) interpretation, subject to confirmation of the source's domain, regularity, and finiteness conventions. The registry labels the problem open; this is supplied status information, not an independently verified assessment.
\subsection{Sources}
\begin{itemize}
\item Local UnsolvedMath registry, permanent ID \(800906\), source problem code \texttt{catalogue entry}, title \emph{Open Problems in Pluripotential Theory --- Question 37}. The supplied record lists Dinew, Guedj, and Zeriahi as proposers, gives the year 2015, and labels the problem \texttt{open}; these are registry attributions and status metadata.
\item Dinew, Guedj, and Zeriahi, \emph{Open Problems in Pluripotential Theory} (2015), Question 37, \url{https://arxiv.org/abs/1511.00705}. The supplied question carries the attribution \textnormal{[Guedj--Zeriahi]}. The provenance records access on 2026-10-16 and extraction by \texttt{arxiv\_latex}; these are supplied provenance notes, not a claim of independent access or verification.
\item Supplied upstream revision: \texttt{d371ee917ac9647fbb06cd6fc529deec6f3dda82}. The registry record was created and updated on 2026-10-16. No separate dated literature assessment or additional literature references are supplied.
\end{itemize}
\end{document}
